\documentclass[10pt,a4paper]{amsart}
\usepackage[margin=19mm]{geometry}
\usepackage{amsmath,amssymb,mathtools}
\usepackage[scr=rsfs,cal=euler]{mathalfa}
\usepackage{xfrac}
\usepackage[unicode,hidelinks,bookmarks=false]{hyperref}
\newcommand{\ie}{i.e.\ }
\newcommand{\cf}{cf.\ }
\newcommand{\resp}{respectively\ }
\newcommand{\Subsection}{\subsection}
\providecommand{\nobreakdash}{\nobreak\hbox{-}}
\setlength{\emergencystretch}{2em}
\multlinegap0pt
\usepackage{amssymb}
\usepackage{bm}
\usepackage{xcolor}
\newtheorem{theorem}{Theorem}
\newcommand{\R}{\mathbb{R}}
\newcommand{\Z}{\mathbb{Z}}
\newcommand{\dd}{\mathrm{d}}
\title{Series decomposition of a class of special integrals}
\author{Xiaolei Yang}
\date{Source: arXiv:2512.12700v2 (CC BY 4.0)}
\begin{document}
\maketitle

\noindent\emph{Source and licence.} Series decomposition of a class of special integrals. Version arXiv:2512.12700v2 (CC BY 4.0), distributed by arXiv under the Creative Commons Attribution 4.0 International licence, http://creativecommons.org/licenses/by/4.0/, as recorded in the arXiv licence metadata for this exact version and shown on the abstract page https://arxiv.org/abs/2512.12700v2. Full licence notice: \texttt{licenses/arxiv-2512.12700-CC-BY-4.0.txt}.\par\smallskip

\noindent\textit{Editorial scope.} Counted unit: the article's theorem thm1 (source0.tex lines 114--124). The article's complete original proof of it is delivered with it (source0.tex lines 126--236, one \texttt{proof} environment of 111 lines). No mathematics is written, repaired or re-derived: the statement and the proof are the article's own lines, in the article's own notation, and no retained line is substituted or re-flowed.  Retained uncounted as the source-local context the proof needs: lines 67--72.  Nothing was omitted from any retained span, so the package carries no omission record.  No retained line contains a citation, so the wrapper carries no bibliography and no bibliography entry.  No retained line contains a figure environment, an image-inclusion command or drawing code, so no image is delivered and the package declares no asset dependency.  Editorial formatting only, in addition: an AMS \texttt{amsart} preamble that restates the article's own declarations and macros used by the retained lines, namely the environments \texttt{theorem} on separate counters, together with the standard packages those lines need.  The retained environments print in this document as Theorem 1 in order of appearance, whereas the article prints the same items with the numbering of its own full document.  The complete native source of the article as distributed in the arXiv e-print for this exact version is bundled unchanged as \texttt{sources/190963/source0.tex}.
\begin{equation}\label{1}
	\begin{cases}
		u_t(t,x) = d(J*u)(t,x) - u(t,x), & t > 0, x \in \mathbb{R} \\
		u(0,x) = u_0(x), & x \in \mathbb{R}
	\end{cases}
\end{equation}

\begin{theorem}\label{thm1}
 Let $u_{0}(x)\in \mathscr{S}'(\mathbb{R})$ and $\hat{u}_{0}(x)\in L^{1}(\R)$. Under the condition of \textbf{(H1)}, denote the solution of Eq.\eqref{1} as \(u(t,x)\), then,
\begin{equation}\label{T1}
	|u(t,x)|\leq \|\hat{u}_0\|_{L^1(\R)}e^{(d-1)t}.
\end{equation}
Furthermore, $e^{(d-1)t}$ is sharp in the sense that for any $\varepsilon>0$, there are $J(x)$ and $  u_0(x)$  such that $u(t,x)$ satisies 
\begin{equation}\label{T2}
		g(x)e^{(d-1-\varepsilon)t}<|u(t,x)|<h(x)e^{(d-1)t},
\end{equation}
 where $g(x), h(x)>0$ are determined only by $u_{0}(x)$. 
\end{theorem}

\begin{proof}
When \( J(x) = \delta(x) \), noting that \( (\delta * u)(t,x) = u(t,x) \), system \eqref{1} then reduces to the local case
\begin{equation*}
\begin{cases}
u_t(t,x) = (d-1)u(t,x), & t > 0, x \in \mathbb{R}, \\
u(0,x) = u_0(x), & x \in \mathbb{R},
\end{cases}
\end{equation*}
which has the solution \( u(t,x) = e^{(d-1)t}u_0(x) \). The rate \(e^{(d-1)t}\) emerges here.

{\bf Step 1:} we show that \(e^{(d-1)t}\) is an upper bound. Let $u_{0}(x)\in \mathscr{S}'(\mathbb{R})$ and $\hat{u}_{0}(x)\in L^{1}(\R)$. Applying the Fourier transform to \eqref{1} gives
\(\hat{u}(t,\xi) = e^{t(d\hat{J}(\xi)-1)}\hat{u_0}(\xi)\in L^1(\R)\). The result follows from the fact that $\parallel \hat{f}\parallel_{L^{\infty}(\mathbb{R})}\leq\parallel f\parallel_{L^{1}(\mathbb{R})}$ for $f\in L^{1}(\R)$.

{\bf Step 2:} We choose suitable functions to ensure the sharpness of \textbf{Theorem \ref{thm1}}. Here we impose the following condition:\\
\textbf{(H2)} \(\hat{u}_0\) is a nonagetive even function and radially decreasing, whatsmore, for any \(x>0\),
\begin{equation}\label{h2}
	\int_0^{\frac{3}{4x}}\hat{u}_0(\xi)\cos(2\pi x\cdot\xi)\dd\xi>0.
\end{equation}

From now on, let \(J\) is a centered normal distribution whose parameter \(\sigma\) to be determined later. Then,
\begin{align*}
	u(t,x) &= \mathscr{F}^{-1}\left(\hat{u}(t,\xi)\right) = \int_{\mathbb{R}} e^{t(d\hat{J}(\xi)-1)}\hat{u}_0(\xi)e^{2\pi i x \xi}\dd \xi \\
	&= \int_{\mathbb{R}} e^{t(d\hat{J}(\xi)-1)}\hat{u}_0(\xi)(\cos 2\pi x \xi + i \sin 2\pi x \xi)\dd \xi \\
	&= \int_{\mathbb{R}} e^{t(d\hat{J}(\xi)-1)}\hat{u}_0(\xi)\cos(2\pi x \xi)\dd \xi.
\end{align*}
where we have used the facts that \(\sin (2\pi x \xi)\) is an odd function and \(e^{t(d\hat{J}(\xi)-1)}\hat{u}_0(\xi)\in L^1(\R)\). It is easy to know that \(u(t,x)\) is even, then, it suffices to consider the case \(x \geq 0\). 

 When \(x > 0\), let
\[
A_k = \int_{\frac{2k-1}{4x}}^{\frac{2k+1}{4x}} \hat{u}_0 (\xi) \cos (2\pi x \xi) \dd \xi, \quad k \in \Z;
\]
\[
B_k = \int_{\frac{2k-1}{4x}}^{\frac{2k+1}{4x}} e^{t(d\hat{J}(\xi)-1)} \hat{u}_0 (\xi) \cos (2\pi x \xi) \dd \xi, \quad k \in \mathbb{Z}.
\]
Here, \(A_k\) and \(B_k\) are functions of \(x\), and the same applies to \(C_k\) mentioned later. Note that \(\hat{J}(\xi)=e^{-2\pi^2 \sigma^2\xi^2}\) is a radially decreasing positive function, and for every \(k\), the integration intervals for both \(A_k\) and \(B_k\) are \(\left( (\frac{2k-1}{4x}),(\frac{2k+1}{4x})\right) \), then \(2\pi x\xi\in (k\pi-\frac{\pi}{2},k\pi+\frac{\pi}{2})\), thus integrands associated with \(A_k\) and \(B_k\) do not change sign in the corresponding interval. With these properties, we can show that:

(i) \( A_k = A_{-k}, B_k = B_{-k} \); Moreover,
\[
A_k, B_k
\begin{cases} 
	> 0, & k = 2m;\\ 
	< 0, & k = 2m + 1.
\end{cases}
\]

(ii) \(B_k\) can be estimated by \(A_k\),
\begin{equation}\label{5}
	|A_k|e^{t\left( d\hat{J}\left( \frac{2k+1}{4x} \right)-1 \right)} < |B_k| < |A_k|e^{t\left( d\hat{J}\left( \frac{2k-1}{4x} \right)-1 \right)}, \quad |k| \geq 1
\end{equation}
\begin{equation}\label{6}
	|A_0|e^{t\left( d\hat{J}\left( \frac{1}{4x} \right)-1 \right)} < |B_0| < |A_0|e^{t(d-1)}.
\end{equation}

(iii) \(|A_k|\) and \(|B_k|\) are both monotonically decreasing for \(k\geq0\). Indeed, when \(k\geq 1\),
\begin{align*}
	|A_{k+1}| &= \left| \int_{\frac{2k+1}{4x}}^{\frac{2k+3}{4x}} \hat{u}_0 (\xi) \cos (2\pi x \xi) \dd \xi\right|=\int_{\frac{2k+1}{4x}}^{\frac{2k+3}{4x}} \left| \hat{u}_0 (\xi) \cos (2\pi x \xi)\right|  \dd \xi  \\
	&=  \int_{\frac{2k-1}{4x}}^{\frac{2k+1}{4x}} \left|\hat{u}_0 \left(\xi+\frac{1}{2x}\right) \cos (2\pi x \xi + \pi)\right| \dd \xi= \int_{\frac{2k-1}{4x}}^{\frac{2k+1}{4x}} \left| \hat{u}_0 \left(\xi+\frac{1}{2x}\right) \cos (2\pi x \xi)\right|  \dd \xi  \\
	&< \int_{\frac{2k-1}{4x}}^{\frac{2k+1}{4x}} \left| \hat{u}_0 \left(\xi\right) \cos (2\pi x \xi)\right|\dd \xi=|A_k|.
\end{align*}
Besides, \(|A_0|>|A_1|\) also holds since \(|\xi+\frac{1}{2x}|<|\xi|\) for \(\xi\in (-\frac{1}{4x},\frac{1}{4x})\). Taking into consideration that \( e^{t(d\hat{J}(\xi)-1)}\hat{u}_0(\xi)\) also decreases radially, we obtain the monotonicity of \(|B_k|\) for \(k\geq0\).

(iv) Since \(\hat{u}_0, e^{t(d\hat{J}-1)} \hat{u}_0\in L^1(\R)\), then, series \(A_n\) and \(B_n\) are absolutely convergent. It follows that \(A_k, B_k\rightarrow0\) as \(k\rightarrow \infty\).


Now, define sequence \(\{C_n\}_{n=0}^\infty\) as
\[C_n = 
\begin{cases} 
	B_0, & n = 0, \\ 
	B_n+B_{-n}, & n > 0. 
\end{cases}\]
Then, we have
\begin{align*}
u(t,x) = \int_{\mathbb{R}} e^{t(d\hat{J}(\xi)-1)}\hat{u}_0(\xi)\cos(2\pi x \xi)\dd \xi
= \sum_{k=-\infty}^{\infty} B_k=\sum_{k=0}^{\infty} C_k.
\end{align*}
 Combined with E.q. \eqref{5}, \eqref{6}, we know that 
 \[0<A_0e^{t\left( d\hat{J}\left( \frac{1}{4x} \right)-1 \right)} < C_0=B_0 < A_0e^{t(d-1)},\]
 	\[A_1 e^{t\left( d\hat{J}\left( \frac{1}{4x} \right)-1 \right)} < B_1=B_{-1} < A_1 e^{t\left( d\hat{J}\left( \frac{3}{4x} \right)-1 \right)}, \]
 Thus \(C_1=B_1+B_{-1}\geq2A_1 e^{t\left( d\hat{J}\left( \frac{1}{4x} \right)-1 \right)}\). Note that condition \textbf{(H2)} means \(\frac{1}{2}A_0+A_1>0\), hence \(\{C_n\}_{n=0}^{\infty}\) is also an alternating series that Leibniz test can be applied. What's more,
  \begin{equation}\label{7}
  	(A_0 + 2A_1)e^{t(d\hat{J}(\frac{1}{4x})-1)} <C_0+C_1< u(t,x) <C_0< A_0e^{t(d-1)}.
  \end{equation}
  
 Fix \(x_0>0\) small enough such that
 \begin{equation*}
 	\int_{-\frac{1}{8x_0}}^{\frac{1}{8x_0}}\hat{u}_0(\xi)\dd \xi\geq \frac{3\|\hat{u}_0\|_{L^1(\R)}}{4},
 \end{equation*}
 Then, for \(0\leq x<x_0 \),
 \begin{align*}
 	u(t,x) &= \int_{\mathbb{R}} e^{t(d\hat{J}(\xi)-1)}\hat{u}_0(\xi)\cos(2\pi x \xi)\dd \xi \\
 	&= \int_{|\xi|\leq \frac{1}{8x_0}} e^{t(d\hat{J}(\xi)-1)}\hat{u}_0(\xi)\cos(2\pi x \xi)\dd \xi+ \int_{|\xi|> \frac{1}{8x_0}} e^{t(d\hat{J}(\xi)-1)}\hat{u}_0(\xi)\cos(2\pi x \xi)\dd \xi\\
 	&\geq \int_{|\xi|\leq \frac{1}{8x_0}} e^{t(d\hat{J}(\xi)-1)}\hat{u}_0(\xi)\cos(2\pi x \xi)\dd \xi-\int_{|\xi|> \frac{1}{8x_0}}\left|  e^{t(d\hat{J}(\xi)-1)}\hat{u}_0(\xi)\cos(2\pi x \xi)\right| \dd \xi\\
 	&\geq \frac{\sqrt{2}}{2}\int_{|\xi|\leq \frac{1}{8x_0}} e^{t(d\hat{J}(\xi)-1)}\hat{u}_0(\xi)\dd \xi-\int_{|\xi|> \frac{1}{8x_0}}\left|  e^{t(d\hat{J}(\xi)-1)}\hat{u}_0(\xi)\cos(2\pi x \xi)\right| \dd \xi\\
 	&\geq \frac{\sqrt{2}}{2}e^{t(d\hat{J}(\frac{1}{8x_0})-1)}\int_{|\xi|\leq \frac{1}{8x_0}} \hat{u}_0(\xi)\dd \xi-e^{t(d\hat{J}(\frac{1}{8x_0})-1)}\int_{|\xi|> \frac{1}{8x_0}}\left| \hat{u}_0(\xi)\cos(2\pi x \xi)\right| \dd \xi\\
 	&\geq \frac{3\sqrt{2}}{8}e^{t(d\hat{J}(\frac{1}{8x_0})-1)} \|\hat{u}_0\|_{L^1(\R)}-e^{t(d\hat{J}(\frac{1}{8x_0})-1)}\int_{|\xi|> \frac{1}{8x_0}}\left| \hat{u}_0(\xi)\right| \dd \xi\\
 	&\geq \left( \frac{3\sqrt{2}}{8}-\frac{1}{4}\right)\|\hat{u}_0\|_{L^1(\R)}e^{t(d\hat{J}(\frac{1}{8x_0})-1)} .
 \end{align*}
 
 \eqref{7} shows that for \(x\geq x_0\), 
 \[u(t,x)\geq (A_0 + 2A_1)e^{t(d\hat{J}(\frac{1}{4x_0})-1)}.\] 
 In summary, we have
 \[u(t,x)\geq c(x)e^{t(d\hat{J}(\frac{1}{4x_0})-1)}\]
 where
 \[c(x)  :=  \left\{ \begin{array}{l}
 	\frac{{3\sqrt 2-2 }}{8},\quad 0 \le x < {x_0},\\
 	{A_0} + 2{A_1},\quad x > {x_0}.
 \end{array} \right.\]
Remember that \(\hat{J}(\xi)=e^{-2\pi^2 \sigma^2\xi^2}\), then, for any \(\varepsilon > 0\), ther exists \(\sigma > 0\) such that \(\hat{J}\left(\frac{1}{4x}\right) > 1 - \frac{\varepsilon}{d}\), then
 \[e^{t(d\hat{J}(\frac{1}{4x})-1)} >e^{t(d-\varepsilon-1)},\]
  that is, \(e^{t(d-1)}\) is sharp.
\end{proof} 

\end{document}
